% 小行星（综述）
% license CCBYSA3
% type Wiki

本文根据 CC-BY-SA 协议转载翻译自维基百科\href{https://en.wikipedia.org/wiki/Asteroid}{相关文章}。

\begin{figure}[ht]
\centering
\includegraphics[width=8cm]{./figures/9fdf35f418301f87.png}
\caption{被探测访问过的小行星影像展示其差异：（上排）433 Eros 与 243 Ida 及其卫星 Dactyl，（下排）Ceres 与 101955 Bennu。尺寸不按比例绘制。} \label{fig_Astero_1}
\end{figure}
小行星是一类小行星体（minor planet），其尺寸大于流星体（meteoroid），但既不是行星，也不是已确认为彗星的天体。它们位于太阳系内部轨道运行，或与木星共享轨道（如特洛伊小行星）。小行星由岩石、金属或冰质物质构成，不具备大气层，并可大致分类为 C 型（碳质）、M 型（金属）、S 型（硅酸盐）。小行星的大小与形状差异极大，从直径不足一千米的松散碎石堆，到直径接近 1000 千米、被归类为矮行星的 Ceres。若天体在受到太阳辐射加热时产生彗发（尾），则被归类为彗星而非小行星，尽管近期观测显示这两类天体之间可能存在连续光谱。\(^\text{[1][2]}\)

在约一百万颗已知小行星中，\(^\text{[3]}\)数量最多的位于火星与木星轨道之间、距离太阳约 2 至 4 AU 的区域，称为主小行星带（main asteroid belt）。所有小行星总质量仅约为地球月球质量的 3\%。主带中大多数小行星沿轻微椭圆、稳定的轨道运行，方向与地球一致，绕太阳一周约需 3 至 6 年。\(^\text{[4]}\)

小行星最早通过地面观测方式进行研究，首次近距离观测来自伽利略号（Galileo）探测器。其后，美国 NASA 与日本 JAXA 相继发射了多次专门的小行星探测任务，并持续规划新的任务。NASA 的 NEAR Shoemaker 研究了 Eros，Dawn 探测了 Vesta 与 Ceres；JAXA 的隼鸟号（Hayabusa）与隼鸟 2 号分别研究并回收了 Itokawa 与 Ryugu 的样本；OSIRIS{-}REx 探测 Bennu，并于 2020 年采样，2023 年将样本送回地球。NASA 的 Lucy（2021 年发射）将研究十颗不同小行星，其中两颗来自主带、八颗为木星特洛伊小行星。Psyche 探测器于 2023 年 10 月发射，目标是金属小行星 Psyche。欧洲航天局的 Hera 于 2024 年 10 月发射，旨在研究 DART 碰撞后的结果。中国国家航天局的天问二号于 2025 年 5 月发射，\(^\text{[5]}\)其任务为探测共轨近地小行星 469219 Kamoʻoalewa 与活跃小行星 311P/PanSTARRS，并对 Kamoʻoalewa 的风化层（regolith）进行采样。\(^\text{[6]}\)

近地小行星若与地球发生撞击，可能造成灾难性影响，其中最著名的案例是 Chicxulub 撞击事件，被普遍认为导致了白垩纪–古近纪生物大灭绝。为应对这一潜在威胁，作为实验性防御措施，2022 年 9 月，双小行星偏转测试（Double Asteroid Redirection Test, DART）探测器通过撞击成功改变了无威胁小行星 Dimorphos 的轨道。
\subsection{术语}
\begin{figure}[ht]
\centering
\includegraphics[width=8cm]{./figures/795a52138582bb97.png}
\caption{一幅按相同尺度制作的组合图像，展示了 2012 年前以高分辨率成像的小行星。从大到小依次为：4 Vesta、21 Lutetia、253 Mathilde、243 Ida 及其卫星 Dactyl、433 Eros、951 Gaspra、2867 Šteins、25143 Itokawa。} \label{fig_Astero_2}
\end{figure}
\begin{figure}[ht]
\centering
\includegraphics[width=8cm]{./figures/1f1fc84d80338171.png}
\caption{按相同比例展示的 Vesta（左）、Ceres（中）与月球（右）} \label{fig_Astero_3}
\end{figure}
2006 年，国际天文学联合会（IAU）引入了当前更为推荐的宽泛术语“太阳系小天体”（small Solar System body），将其定义为太阳系中既不是行星、也不是矮行星、亦不是天然卫星的天体；其中包括小行星、彗星以及更近期发现的其他类型天体\(^\text{[7]}\)。根据 IAU 表述，“术语『小行星（minor planet）』仍可继续使用，但一般情况下将优先采用『太阳系小天体（Small Solar System Body）』”\(^\text{[8]}\)。

在历史上，第一颗被发现的小行星——Ceres——最初曾被视为一颗新行星\(^\text{[a]}\)。随后，天文学家又陆续发现了其他相似天体。在当时的观测设备条件下，它们呈现为类似恒星的光点，几乎看不到行星盘结构，但由于具有明显的视运动，又可与恒星区分开来。这促使天文学家威廉·赫歇尔（Sir William Herschel）提出了“asteroid”一词\(^\text{[b]}\)，该词源自希腊语 ἀστεροειδής（asteroeidēs），意为“类星”“似星状”，并由古希腊语 ἀστήρ（astēr，“星、行星”）派生。在 19 世纪后半叶初期，“asteroid”与“planet”（有时并未特别限定为“minor”）仍被交替使用\(^\text{[c]}\)。

传统上，围绕太阳运行的小型天体被划分为彗星、小行星或流星体，其中直径小于一米的天体被称为流星体。术语“asteroid”从未被正式定义\(^\text{[13]}\)，但在非正式使用中可指“围绕太阳运行、形状不规则、呈岩质组成且不符合 IAU 定义中行星或矮行星条件的天体”\(^\text{[14]}\)。小行星与彗星的主要区别在于，彗星因近表层冰在太阳辐射作用下发生升华而形成彗发（尾）。少数天体最初被归类为小行星，但后续观测显示存在彗状活动。相反，一些（或许全部）彗星最终会耗尽其表面挥发性冰，从而演化为类小行星天体。另一区分方式是轨道偏心率：彗星通常具有比大多数小行星更高的轨道偏心率；而高度偏心的小行星极可能为休眠或灭绝的彗星\(^\text{[15]}\)。

位于木星轨道之外的小行星体有时亦被称为“asteroids”，尤其是在大众化科普材料中\(^\text{[d]}\)。然而，将“asteroid”一词限定用于太阳系内部的小行星体的做法正在变得愈发普遍\(^\text{[17]}\)。因此，本文主要将讨论范围限定于传统意义上的小行星：即小行星带天体、木星特洛伊小行星以及近地天体。

自 1801 年 Ceres 被发现后的近两个世纪中，所有已知小行星在其轨道的大部分时间里都位于木星轨道以内，尽管极少数如 944 Hidalgo 会在轨道部分区段上延伸至更远处。自 1977 年 2060 Chiron 被发现起，天文学家开始发现永久位于木星轨道之外的小型天体，即后称为半人马（centaurs）。1992 年，又发现了 15760 Albion，这是首个已知位于海王星轨道之外的天体（除 Pluto 外）；随后人们观测到了大量类似天体，即现称为跨海王星天体（trans-Neptunian object）。更外层依次还有柯伊伯带天体、散射盘天体以及距离更为遥远的奥尔特云——其被假设为休眠彗星的主要储库。它们位于太阳系寒冷的外缘区域，冰质物质能够长期保持固态，因此彗状活动极为微弱；若半人马或跨海王星天体接近太阳，其挥发性冰质材料将发生升华，按传统分类方法将被归类为彗星。

之所以将柯伊伯带天体称为“objects”，部分原因是为了避免必须将其归为小行星或彗星\(^\text{[1]}\)。它们被认为在成分上大体类似于彗星，但部分天体可能与小行星更相近\(^\text{[18]}\)。绝大多数此类天体并不具有典型彗星所表现的高偏心率轨道，且目前发现者均比传统彗核更大。近年来的观测结果，如 Stardust 探测器采集的彗星尘埃样本分析，正不断模糊小行星与彗星之间的界限\(^\text{[1]}\)，越来越多研究表明二者之间可能形成连续谱，而非清晰的分类边界\(^\text{[2]}\)。

2006 年，IAU 为最大的小行星体引入了“矮行星”这一类别——即因自身引力而形成近似椭球形状的天体。迄今为止，仅小行星带中体积最大的天体 Ceres（直径约 975 km，约 606 mi）被划入此类\(^\text{[19][20]}\)。
\subsection{观测历史}
尽管小行星数量极为庞大，但其发现在人类天文学史上相对较晚，第一颗小行星——Ceres——直到 1801 年才被识别出来\(^\text{[21]}\)。在暗夜且观测条件良好的情况下，只有一颗小行星，即表面反照率较高的 4 Vesta，通常能够以肉眼看见。极少数情况下，当小型小行星非常接近地球时，也可能在短暂时间内被肉眼观测到\(^\text{[22]}\)。截至 2025 年 5 月，小行星中心（Minor Planet Center）已记录太阳系内外共计 1,460,356 颗小行星的数据，其中约有 826,864 颗具备足够观测信息而获得编号编号\(^\text{[23]}\)。
\subsubsection{Ceres 的发现}
\begin{figure}[ht]
\centering
\includegraphics[width=6cm]{./figures/63dcf87fd16a81c7.png}
\caption{《Giuseppe Velasquez 与 Giuseppe Piazzi 及 Ceres》，油画，约作于 1803 年} \label{fig_Astero_4}
\end{figure}
1772 年，德国天文学家 Johann Elert Bode（援引 Johann Daniel Titius 的观点）发表了一组数列进程，即如今已被否定的 Titius–Bode 定律。除了火星与木星之间存在无法解释的空隙外，该公式似乎能够预测当时已知行星的轨道\(^\text{[24][25]}\)。Bode 对这一“缺失行星”的存在给出了如下解释：

“这一点尤其来源于一项令人惊异的关系，即已知六颗行星在其距日距离上所遵循的规律。若取土星到太阳的距离为 100，则水星距太阳为 4；金星为 4 + 3 = 7；地球为 4 + 6 = 10；火星为 4 + 12 = 16。此处出现一个在规则进程中的空档：在火星之后应为 4 + 24 = 28 的位置，但迄今未见任何行星。难道可以相信宇宙的创造者故意让这一区域保持空缺吗？显然不能。从此处继续，木星距离为 4 + 48 = 52，最终土星为 4 + 96 = 100。”\(^\text{[26]}\)

依据该公式，Bode 预测将在距离太阳约 2.8 天文单位（AU），即约 4.2 亿千米的位置发现另一颗行星\(^\text{[25]}\)。随着 William Herschel 在接近公式预测位置发现了土星外的新行星天王星，Titius–Bode 定律一度得到了强有力支持\(^\text{[24]}\)。1800 年，由德国天文学刊物《Monatliche Correspondenz》（《每月通信》）主编 Franz Xaver von Zach 牵头，他向 24 名资深天文学家发出邀请（并将其称为“天国警察”）\(^\text{[25]}\)，希望他们协同开展针对该“预期行星”的系统性搜寻\(^\text{[25]}\)。虽然他们未能发现 Ceres，但最终陆续发现了编号为 2、3、4 的三颗小行星：2 Pallas、3 Juno 与 4 Vesta\(^\text{[25]}\)。

被选入搜寻小组的天文学家之一是 Giuseppe Piazzi，他是西西里巴勒莫学院的一位天主教神父。在收到加入该小组的正式邀请之前，Piazzi 已于 1801 年 1 月 1 日发现了 Ceres\(^\text{[27]}\)。当时他正在寻找“拉卡伊先生（Mr la Caille）《黄道星目录》中第 87 颗星”\(^\text{[24]}\)，却注意到“其前方又出现了一颗天体”\(^\text{[24]}\)。Piazzi 实际上发现了一颗会移动的、类星状的天体，并最初认为其是一颗彗星\(^\text{[28]}\)：

“其光度略显微弱，颜色类似木星，与通常被视为八等星的许多恒星非常相近。因此，我毫不怀疑其应为一颗普通恒星。[……]然而在第三天夜晚，我的怀疑已转为确信——它绝非一颗固定恒星。不过在公开之前，我仍决定等待至第四天傍晚，当时我欣喜地看到它继续以前几日相同的速率移动。”\(^\text{[24]}\)

Piazzi 共计观测 Ceres 24 次，最后一次记录于 1801 年 2 月 11 日，之后因病中断研究。他于 1801 年 1 月 24 日向两位同行天文学家致信宣布此发现：米兰的同乡 Barnaba Oriani 与柏林的 Bode\(^\text{[21]}\)。在信中，他将其报告为一颗彗星，但同时表示：“鉴于其运动如此缓慢且相当规则，我已多次想到它可能比彗星更为特殊。”\(^\text{[24]}\)

4 月，Piazzi 将完整观测数据寄给 Oriani、Bode 与法国天文学家 Jérôme Lalande，并于当年 9 月发表于《Monatliche Correspondenz》期刊\(^\text{[28]}\)。

其时，由于地球绕太阳公转引起的视位置变化，Ceres 的位置已发生偏移，并过于靠近太阳炫光，其他天文学家无法立即证实 Piazzi 的观测结果。年末 Ceres 理论上应再次可见，但由于间隔时间长、轨道参数尚不精确，其位置难以预测。为重新寻获 Ceres，时年仅 24 岁的数学家 Carl Friedrich Gauss 开发出一种高效的轨道确定方法\(^\text{[28]}\)。仅数周后，他完成预测并将结果寄给 von Zach。1801 年 12 月 31 日，von Zach 与另一位“天国警察”Heinrich W. M. Olbers 在几乎完全符合预测的位置成功重新发现 Ceres\(^\text{[28]}\)。Ceres 位于距太阳 2.8 AU 的轨道位置，与 Titius–Bode 定律几乎完全吻合；然而海王星在 1846 年被发现后，其距离预测值偏差高达 8 AU，促使大多数天文学家最终将该定律视为巧合\(^\text{[29]}\)。Piazzi 将新发现的天体命名为 Ceres Ferdinandea，以“纪念西西里守护女神与波旁王朝国王 Ferdinand”\(^\text{[26]}\)。
\subsubsection{进一步搜寻}
\begin{figure}[ht]
\centering
\includegraphics[width=8cm]{./figures/321bbb2080ca89a0.png}
\caption{首十颗被发现的小行星的尺寸，与月球的比较} \label{fig_Astero_5}
\end{figure}
在随后的数年中，von Zach 小组又陆续发现了另外三颗小行星，即 2 Pallas、3 Juno 与 4 Vesta，其中 Vesta 于 1807 年被发现\(^\text{[25]}\)。此后新的小行星发现陷入停滞，直到 1845 年才出现进展。业余天文学家 Karl Ludwig Hencke 自 1830 年起开始搜寻新小行星，十五年后在寻找 Vesta 的过程中意外发现了后来被命名为 5 Astraea 的小行星，这是 38 年以来首次获得的新发现。Carl Friedrich Gauss 受邀为其命名。此事件促使更多天文学家加入研究，到 1851 年底共计发现 15 颗小行星。1868 年，当 James Craig Watson 发现第 100 颗小行星时，法国科学院在一枚纪念章上雕刻了三位当时最成功的小行星发现者 Karl Theodor Robert Luther、John Russell Hind 与 Hermann Goldschmidt 的肖像，以纪念这一成果\(^\text{[30]}\)。

1891 年，Max Wolf 首创使用天体摄影法（astrophotography）来探测小行星，长时间曝光的照相底片上，小行星会呈现为短条状轨迹\(^\text{[30]}\)。与以往的目视观测相比，这一技术显著提高了发现效率：Wolf 自 323 Brucia 起共发现 248 颗小行星\(^\text{[31]}\)，而此前历史上总共也不过发现略多于 300 颗。尽管天文学界已知小行星数量远不止于此，但多数天文学家对其兴趣不大，甚至有人称其为“天空的害虫”（vermin of the skies）\(^\text{[32]}\)，此说法被分别归因于 Eduard Suess\(^\text{[33]}\)与 Edmund Weiss\(^\text{[34]}\)。即便在一个世纪后的 20 世纪，仅有数千颗小行星获得识别、编号与命名。
\subsubsection{19 与 20 世纪}
\begin{figure}[ht]
\centering
\includegraphics[width=8cm]{./figures/dc15219005826826.png}
\caption{此处以雷达成像展示的 2013 EC，仍仅具有临时编号} \label{fig_Astero_6}
\end{figure}
过去，小行星的发现通常遵循四个步骤。首先，使用广视场望远镜或天体照相仪对天空某一区域进行拍摄，同一区域需拍摄成对照片，两张照片的拍摄时间间隔通常为一小时，多天内可进行多次成对拍摄。其次，将同一区域的两张底片或胶片置于体视镜（stereoscope）下进行比对，由于绕日运行天体会在两次拍摄期间产生微小位移，其影像在体视镜观察下会呈现为略高于恒星背景的悬浮效果。第三，在确认目标为移动天体后，需使用数字化显微测量仪对其位置进行精确测量，并相对于已知恒星的坐标进行标定\(^\text{[35]}\)。

这前三个步骤尚不能构成对小行星的正式发现：此时观测者仅获得了一次“视象”（apparition），并会被赋予一个临时编号（provisional designation），其格式由发现年份、发现所在半月区间的字母以及按顺序分配的字母与数字组成（示例：1998 FJ74）。最后一步是将观测位置与观测时间提交至小行星中心（Minor Planet Center），由其计算机程序判断该视象是否与以往记录关联并构成同一轨道。若条件满足，该天体将获得正式星表编号，而首次具备可计算轨道视象的观测者将被确认为发现者，并在国际天文学联合会批准下获得命名权\(^\text{[36]}\)。
\subsection{命名}
到 1851 年为止，英国皇家天文学会认为小行星的发现速度过快，需要采用新的分类或命名体系。1852 年，当 de Gasparis 发现第 20 颗小行星时，Benjamin Valz 为其赋予名称和序号，即 20 Massalia，以示其在小行星发现序列中的位置。有些小行星在发现后未能再次被观测到，因此自 1892 年起，新发现的小行星改以发现年份加大写字母的方式记录，用以标识当年同类天体的登记顺序。例如，1892 年发现的前两颗小行星被标记为 1892A 与 1892B。然而，到 1893 年时，字母数已不足以标识全部发现天体，于是在 1893Z 之后使用 1893AA。此后又尝试过多种变体，包括 1914 年起采用“年份 + 希腊字母”的方式。最终，于 1925 年确立了简明的按时间顺序编号体系\(^\text{[30][37]}\)。

目前，所有新发现的小行星均会先获得临时编号（provisional designation），例如 2002 AT4。其格式由发现年份和表示发现半月周期及序号的字母数字组合构成。当小行星轨道获得确认后，将赋予其正式编号，并可在之后获得名称（例如 433 Eros）。正式命名规约使用括号包围编号，如 (433) Eros，但省略括号的写法也相当普遍。在非正式语境中，人们常省略编号；或在连续文本中仅于首次出现时标注编号\(^\text{[38]}\)。此外，小行星的命名可以由发现者按照国际天文学联合会制定的规范提出\(^\text{[39]}\)。
\subsubsection{符号}  
最早被发现的小行星被赋予类似行星符号的标志性象征符号。到 1852 年时，已存在约二十余种小行星符号，且常出现多个变体\(^\text{[40]}\)。

1851 年，在第 15 颗小行星 Eunomia 被发现之后，Johann Franz Encke 在将出版的 1854 年版《柏林天文年鉴》（Berliner Astronomisches Jahrbuch, BAJ）中进行了重大修改。他引入了传统用于表示恒星的圆盘符号（即圆圈）作为小行星的通用符号，并以发现顺序在圆圈中标注数字，以表示特定小行星。该编号圆圈体系迅速被天文学界采用，其后的第 16 颗小行星（1852 年发现的 16 Psyche）成为首个在发现时即采用该体系标示的小行星。然而 Psyche 同时仍被赋予传统象征符号，此后若干小行星亦然。20 Massalia 是首个未获象征符号的小行星，而最后一个被赋予象征符号的小行星是 1855 年发现的 37 Fides\(^\text{[e][41]}\)。
\subsection{形成} 
许多小行星是原行星（planetesimals）的破碎残骸，这些天体曾存在于年轻太阳的星云中，但未能成长至行星尺度\(^\text{[42]}\)。普遍认为，小行星带中的原行星最初与太阳星云中其他天体一样经历演化，直到木星增至接近其当前质量时，由其轨道共振所产生的摄动将带中超过 99\% 的原行星弹出。数值模拟结果，以及自转速率与光谱性质的不连续性均表明：直径大于约 120 km（75 mi）的小行星可能在早期阶段经吸积形成，而更小的天体则是木星扰动期间或其后，小行星之间碰撞所产生的碎片\(^\text{[43]}\)。Ceres 与 Vesta 的体量足够大，因此经历了熔融与分化，重金属元素下沉形成核心，而岩质矿物则构成外壳\(^\text{[44]}\)。

根据 Nice 模型，许多柯伊伯带天体被捕获进入外部小行星带，位于距离太阳大于 2.6 AU 的区域。其中大多数随后被木星再次弹出，但残留者可能构成 D 型小行星，并且可能包括 Ceres 本身\(^\text{[45]}\)。
\subsection{太阳系内的分布}
\begin{figure}[ht]
\centering
\includegraphics[width=6cm]{./figures/e8184c45bcc8665a.png}
\caption{内太阳系中小行星族群位置的俯视示意图} \label{fig_Astero_7}
\end{figure}
在太阳系内部轨道上已发现多类动力学性质不同的小行星族群。它们的轨道受太阳系中其他天体的引力扰动以及雅可夫斯基效应（Yarkovsky effect）的影响。主要族群包括：
\subsubsection{小行星带}
\begin{figure}[ht]
\centering
\includegraphics[width=6cm]{./figures/a627147494f18831.png}
\caption{一幅展示内太阳系行星及小行星族群分布的示意图。距日距离按比例绘制，天体尺寸不按比例。} \label{fig_Astero_8}
\end{figure} 
已知多数小行星位于火星与木星轨道之间的小行星带内，其轨道通常具有较低偏心率（即轨道形状并不显著拉长）。据估计，该区域包含直径大于 1 km（0.6 mi）的小行星约 1.1 至 1.9 百万颗\(^\text{[46]}\)，以及更多体积更小的成员。这些小行星可能是原行星盘的残余物，而在太阳系形成早期，由于木星的强大引力扰动，该区域内的原行星无法通过吸积进一步成长为行星。

与大众印象相反，小行星带大部分区域实际上是空旷的。小行星在极大体积空间内分布稀疏，因此若不进行精确导航，要主动抵达某一目标小行星几乎不可能。尽管如此，当前已知小行星数量达数十万颗，而其总数可能达到数百万甚至更多，这取决于所设定的尺寸下限。目前已知超过 200 颗小行星的直径超过 100 km\(^\text{[47]}\)，且红外波段的观测表明，小行星带内直径大于 1 km 的小行星数量介于 700,000 至 1.7 百万之间\(^\text{[48]}\)。已知小行星的绝对星等多数介于 11 至 19 之间，中位数约为 16\(^\text{[49]}\)。

小行星带的总质量估计约为 \(2.39\times10^{21}\,\mathrm{kg}\)，仅相当于月球质量的约 3\%；相比之下，柯伊伯带与散射盘（Scattered Disk）的总质量超过其 100 倍\(^\text{[50]}\)。在小行星带总质量中，四大天体——Ceres、Vesta、Pallas 与 Hygiea——共占约 62\%，其中 Ceres 单独便占约 39\%。
\subsubsection{特洛伊小行星}  
特洛伊小行星是指与较大行星或卫星共享轨道、但不会与之发生碰撞的天体，它们位于轨道上的两个拉格朗日稳定点 L4 与 L5 处，分别位于该较大天体前方与后方约 \(60^\circ\) 的位置。

在太阳系中，已知的特洛伊小行星绝大多数与木星共享轨道。依据位置不同，位于 L4（木星前方）的族群被称为“希腊营”，位于 L5（木星后方）的族群被称为“特洛伊营”。预计直径大于 1 km 的木星特洛伊小行星数量超过一百万颗\(^\text{[51]}\)，其中已有超过 7{,}000 颗被编目记录。除木星外，其余行星轨道上迄今仅发现九颗火星特洛伊小行星、28 颗海王星特洛伊小行星、两颗天王星特洛伊小行星与两颗地球特洛伊小行星；另有一颗暂时处于特洛伊轨道的金星天体被确认。数值轨道动力学稳定性模拟表明，土星与天王星极可能不存在原初特洛伊天体\(^\text{[52]}\)。
\subsubsection{近地小行星}
近地小行星（Near-Earth asteroids，NEAs）是指轨道接近地球轨道的小行星。若小行星轨道实际与地球轨道发生交叉，则称为地球交叉小行星（Earth-crossers）。截至 2022 年 4 月，已知近地小行星总数为 28{,}772 颗，其中 878 颗的直径不小于 1 km\(^\text{[53]}\)。

少数 NEAs 是已失去表面挥发物的“灭绝彗星”，虽然某些天体存在微弱或间歇性的类彗尾结构，但这并不足以将其正式归类为近地彗星，因此两者的分类边界并不十分明确。其余 NEAs 则主要因受到木星引力扰动而从小行星带中被驱逐出来\(^\text{[54][55]}\)。

许多小行星拥有天然卫星（即小行星卫星）。截至 2021 年 10 月，已知至少 85 颗 NEAs 拥有至少一颗卫星，其中三颗拥有两颗卫星\(^\text{[56]}\)。直径约 4.5 km（2.8 mi）的大型潜在威胁小行星 3122 Florence 具有两颗直径约 100–300 m（330–980 ft）的卫星，它们是在 2017 年该天体接近地球期间，以雷达成像方式被发现\(^\text{[57]}\)。

近地小行星依其轨道半长轴 \(a\)、近日点距离 \(q\) 与远日点距离 \(Q\) 分为以下几类\(^\text{[58][54]}\)：
\begin{itemize}
\item Atira（或 Apohele）族：轨道完全位于地球轨道以内，即远日点距离满足\( Q < 0.983\,\mathrm{AU}\),其中 \(0.983\,\mathrm{AU}\) 为地球近日点距离，这也意味着\(a < 0.983\,\mathrm{AU}\).
\item Aten 族：轨道半长轴小于 1 AU 且轨道与地球轨道交叉，即\(a < 1.0\,\mathrm{AU},\quad Q > 0.983\,\mathrm{AU}\).
\item Apollo 族：轨道半长轴大于 1 AU 且轨道与地球轨道交叉，即\(a > 1.0\,\mathrm{AU},\quad q < 1.017\,\mathrm{AU}\).其中 \(1.017\,\mathrm{AU}\) 为地球远日点距离。
\item Amor 族：轨道完全位于地球轨道之外但与地球轨道邻近，即近日点距离满足\(1.017\,\mathrm{AU} < q < 1.3\,\mathrm{AU}\),这同样意味着\(a > 1.017\,\mathrm{AU}\).部分 Amor 小行星轨道会穿越火星轨道。
\end{itemize}
\subsubsection{火星卫星}
\begin{figure}[ht]
\centering
\includegraphics[width=8cm]{./figures/5d9259af6c07531c.png}
\caption{} \label{fig_Astero_9}
\end{figure}
目前尚不明确火星卫星 Phobos 与 Deimos 是否为被俘获的小行星，或是由火星表面的大型撞击事件形成\(^\text{[60]:1284–1285}\)。二者在诸多性质上与碳质 C 型小行星高度相似，其光谱、反照率及密度均与 C 型或 D 型小行星非常接近\(^\text{[61]}\)。基于此类相似性，一种假说认为这两颗卫星可能是被俘获的主带小行星\(^\text{[62][63]}\)。然而，二者的轨道几乎完全呈圆形，且位于火星赤道面内，这意味着若为俘获起源，则需有机制将最初高度偏心的轨道圆化并调整轨道倾角，使其贴近赤道平面，此过程可能涉及大气阻力与潮汐力的共同作用\(^\text{[64]}\)，但目前尚不确定对于 Deimos 而言是否具备足够时间完成这一过程\(^\text{[60]:1300–1301}\)。此外，小行星俘获过程必须伴随能量耗散，而现今火星大气过于稀薄，不足以通过大气制动俘获与 Phobos 尺度相当的天体\(^\text{[60]}\)。Geoffrey A. Landis 指出，若原始天体为双小行星体系，并在潮汐作用下发生分离，则俘获仍有可能发生\(^\text{[63][65]}\)。

另一项假说认为，Phobos 可能属于太阳系的“第二代”天体，即在火星形成后于轨道上重新凝聚形成，而非与火星共同起源于同一分子云\(^\text{[66]}\)。

还有研究提出，火星过去可能曾被大量与 Phobos 与 Deimos 尺度相当的天体包围，这些天体可能因大型原行星的撞击而被抛入火星轨道\(^\text{[67]}\)。Phobos 的内部具有高度孔隙结构（依据其密度 \(1.88\,\mathrm{g/cm^3}\) 推算，其体积约有 25–35\% 为空隙），这一点与典型小行星起源不相符\(^\text{[68]}\)。对 Phobos 的热红外观测显示其成分主要包含层状硅酸盐（phyllosilicates），该类矿物广泛存在于火星表面。其光谱与所有类球粒陨石（chondrite）类型均不匹配，再次与小行星起源相矛盾\(^\text{[69]}\)。上述两项证据均支持另一种假说：Phobos 源自火星表面受撞击抛射出的物质，在火星轨道再吸积形成\(^\text{[70]}\)，这与地球月球起源的主流理论模式相似。
\subsection{特性}  
\subsubsection{尺寸分布}
\begin{figure}[ht]
\centering
\includegraphics[width=8cm]{./figures/e17af13d9b4023fc.png}
\caption{内太阳系小行星的尺寸与数量分类示意} \label{fig_Astero_10}
\end{figure}
小行星的尺寸范围极其广泛，从最大者近 \(1000\,\mathrm{km}\)（\(620\,\mathrm{mi}\)）直径，到最小仅约 \(1\,\mathrm{m}\)（\(3.3\,\mathrm{ft}\)）的岩石；小于该尺度的天体则被归类为流星体（meteoroid）\(^\text{[f]}\)。其中最大三颗小行星在许多方面与“微型行星”类似：它们近似球形、具有至少部分分化的内部结构\(^\text{[71]}\)，并被认为是原行星的幸存者。而绝大多数小行星尺寸更小、形状不规则，被视为遭受严重撞击的原行星残骸或更大天体的破碎碎片。

矮行星 Ceres 是迄今最大的已知小行星，其直径约为 \(940\,\mathrm{km}\)（\(580\,\mathrm{mi}\)）。第二与第三大的分别是 4 Vesta 与 2 Pallas，其直径均略大于 \(500\,\mathrm{km}\)（\(300\,\mathrm{mi}\)）。Vesta 是四颗主带中偶尔可被肉眼观测到的小行星中最明亮者\(^\text{[72]}\)。在极少数情况下，近地小行星也可能短暂地无需仪器辅助而可见，典型例子如 99942 Apophis。

位于火星与木星轨道之间的小行星带天体总质量约为 \((2394\pm6)\times10^{18}\,\mathrm{kg}\)，约为月球质量的 \(3.25\%\)。其中 Ceres 单独贡献 \(938\times10^{18}\,\mathrm{kg}\)，约占总质量的 \(40\%\)。若再加入接下来三颗最具质量的小行星：Vesta（\(11\%\)）、Pallas（\(8.5\%\)）与 Hygiea（\(3\text{–}4\%\)），合计比例略超 \(60\%\)。再加上随后七颗质量位列前十的小行星，则总占比可达 \(70\%\)\(^\text{[50]}\)。随着单个天体质量下降，小行星数量迅速增加。

小行星数量随尺寸增大而显著减少。尽管其尺度分布整体呈现幂律分布，但在约 \(5\,\mathrm{km}\) 与 \(100\,\mathrm{km}\) 附近存在明显“突起区间”，即数量高于幂律外推值。直径大于约 \(120\,\mathrm{km}\)（\(75\,\mathrm{mi}\)）的小行星大多为原初天体（保留自吸积时期），而更小的小行星则主要是原初天体碎裂后的产物。主带原始小行星族群数量可能为现今的约 \(200\) 倍\(^\text{[73][74]}\)。

\textbf{最大的小行星}
\begin{figure}[ht]
\centering
\includegraphics[width=8cm]{./figures/618beb54f3bef2e9.png}
\caption{由欧洲南方天文台超大望远镜（VLT）拍摄的、小行星带内 42 个最大天体影像} \label{fig_Astero_11}
\end{figure}
\begin{figure}[ht]
\centering
\includegraphics[width=8cm]{./figures/cf17651041083441.png}
\caption{Eros、Vesta 与 Ceres 的尺寸比较} \label{fig_Astero_12}
\end{figure}
小行星带中体积最大的三颗天体——Ceres、Vesta 与 Pallas——均为完整保存的原行星（protoplanets），在多项物理性质上与行星相似，因此与大多数形状不规则的小行星相比显得非典型。第四大天体 Hygiea 虽呈近似球形，但其内部可能未经历分化\(^\text{[75]}\)，与大多数小行星相同。上述四大天体合计约占小行星带总质量的五分之八，其中仅 Ceres 与 Vesta 即占一半。

Ceres 是唯一具备在自引力作用下形成塑性球状结构的小行星，因此也是唯一被认定为矮行星者\(^\text{[76]}\)。其绝对星等远高于其他小行星，约为 3.32\(^\text{[77]}\)，并可能具备一层表面冰层\(^\text{[78]}\)。与类行星天体相同，Ceres 具有分化的内部结构，包括地壳、地幔与地核\(^\text{[78]}\)。迄今为止，尚未发现源自 Ceres 的陨石样本落入地球\(^\text{[79]}\)。

Vesta 同样具备分化内部结构，但其形成位置位于太阳系霜线内侧，因此缺乏水成分\(^\text{[80][81]}\)，其成分主要为玄武质岩石，含有橄榄石等矿物\(^\text{[82]}\)。除了南极巨型撞击坑 Rheasilvia 外，Vesta 形状呈椭球状。Vesta 是 Vestian 家族及其他 V 型小行星的母体，并且是 HED 类陨石的来源，而 HED 陨石约占地球上所有陨石的 5\%。

Pallas 则因自转轴高度倾斜且几乎“侧向自转”而显得独特，这一点与天王星相似\(^\text{[83]}\)。其成分与 Ceres 类似，富含碳与硅，且可能经历部分内部分化\(^\text{[84]}\)。Pallas 是 Palladian 家族小行星的母体。

Hygiea 是目前已知最大的碳质小行星\(^\text{[85]}\)，且与前三者不同，其轨道相对接近黄道面。Hygiea 是 Hygiean 家族中最大成员及推定母体，但其表面并未发现足以解释该家族形成的大型撞击坑（类似于 Vesta 的情形）。因此推测 Hygiea 在形成该家族的撞击事件中曾几乎完全被破坏，并在损失不足约 2\% 质量后重新聚合。2017–2018 年利用超大望远镜 SPHERE 成像仪开展的观测显示，Hygiea 具有近乎球形形状，此结果与其目前或曾经处于流体静力平衡状态，或曾经历完全破碎后再凝聚的情景一致\(^\text{[86][87]}\)。

大型小行星的内部分化现象可能与其缺乏天然卫星的事实有关，因为主带小行星的卫星通常被认为源自剧烈撞击所造成的碎裂过程，进而形成类似碎石堆（rubble pile）的结构\(^\text{[79]}\)。
\begin{figure}[ht]
\centering
\includegraphics[width=14.25cm]{./figures/41f29545a8983c38.png}
\caption{} \label{fig_Astero_13}
\end{figure}
\subsubsection{自转} 
对小行星带内大型小行星自转速率的测量显示其存在上限。直径大于 \(100\,\mathrm{m}\) 的小行星中，自转周期短于 \(2.2\,\mathrm{h}\) 的情况极为罕见\(^\text{[88]}\)。对于自转速率快于该阈值的天体，其表面惯性力将超过引力，使松散表层物质被抛出。然而，若天体为整体坚固体，则理论上可承受更快的自转速率。因此，这一现象表明：多数直径超过 \(100\,\mathrm{m}\) 的小行星应为碎石堆结构（rubble piles），即由小行星间碰撞后产生的碎片重新聚集形成\(^\text{[89]}\)。
\subsubsection{颜色}
小行星会随着暴露于太空风化（space weathering）而逐渐变暗、变红\(^\text{[90]}\)。然而相关证据显示，其表面颜色的大部分变化过程发生极为迅速，主要集中在最初的十万年内，因此光谱测量对于推定小行星年龄的效用有限\(^\text{[91]}\)。
\subsubsection{表面特征}
\begin{figure}[ht]
\centering
\includegraphics[width=8cm]{./figures/1cd3f712ff79e473.png}
\caption{4 Vesta 的布满撞击坑地形} \label{fig_Astero_14}
\end{figure}
除“前四大”小行星（Ceres、Pallas、Vesta 与 Hygiea）外，其余小行星在外观上大体相似，尽管形状普遍不规则。直径约 \(50\,\mathrm{km}\)（\(31\,\mathrm{mi}\)）的 253 Mathilde 是一颗碎石堆小行星（rubble pile），其表面遍布撞击坑，且部分撞击坑直径可达其半径尺度。地基观测显示，直径约 \(300\,\mathrm{km}\)（\(190\,\mathrm{mi}\)）的 511 Davida —— 继“前四大”之后体积较大者之一 —— 具有类似的棱角外形，暗示其同样受到半径尺度撞击坑的饱和破坏\(^\text{[92]}\)。Mathilde 与 243 Ida 等中等尺寸的小行星经近距探测显示，其表面覆盖深厚风化层（regolith）。在“前四大”中，Pallas 与 Hygiea 的表面特征仍鲜为人知；Vesta 南极具有围绕半径尺度撞击盆地的环状压缩断裂结构，但除此之外整体呈近似球状。

“黎明号”（Dawn）探测器显示，Ceres 表面布满撞击坑，但其大型撞击坑数量少于理论预期[93]。基于当前小行星带形成模型推测，Ceres 原应拥有 10 至 15 个直径超过 \(400\,\mathrm{km}\)（\(250\,\mathrm{mi}\)）的大型撞击盆地[93]，但其已确认的最大撞击盆地——Kerwan Basin——直径仅为 \(284\,\mathrm{km}\)（\(176\,\mathrm{mi}\)）[94]。最可能的解释是：Ceres 表层地壳发生黏性松弛（viscous relaxation），即大型撞击坑会在地质时间尺度上逐渐被抚平\(^\text{[93]}\)。
\subsubsection{成分}
小行星通常根据其特征发射光谱进行分类，大多数可归入三大类：C 型、M 型与 S 型，对应于富含碳的碳质小行星、富含金属的金属小行星与富含硅酸盐的石质小行星。小行星的真实物理成分十分多样，且在多数情况下尚未被充分理解。Ceres 似乎由岩质核心与表层冰状地幔组成；Vesta 则被认为具有镍铁核心、橄榄石地幔与玄武质地壳[95]。人们曾认为体积最大的未分化小行星为 10 Hygiea，其组成似乎整体均一、呈碳质球粒陨石性质，但其也可能是经历全球性撞击破碎后再度重组的“曾经分化过的小行星”。其他一些小行星似乎为原行星的残余岩质或金属富集核心或地幔。大多数小型小行星则被认为是由碎片在弱引力下松散团聚而成的“碎石堆结构”（rubble piles），而最大型小行星则可能为坚固整体结构。此外，一些小行星拥有卫星或构成协轨双星系统；此类碎石堆、卫星系统、双星系统及分散小行星族群被认为源自母体小行星甚至行星在撞击中被破坏的结果[96]。

在主小行星带中，似乎存在两大主成分族群：一类为暗色、富挥发物族群，包括 C 型与 P 型小行星，其反照率低于 0.10、密度低于 \(2.2\,\mathrm{g/cm^3}\)；另一类为致密、贫挥发物族群，包括 S 型与 M 型小行星，反照率高于 0.15、密度高于 \(2.7\,\mathrm{g/cm^3}\)。在这些族群内部，大型小行星往往密度更大，推定这是由于高压压实作用所致。对于质量大于 \(10\times10^{18}\,\mathrm{kg}\) 的小行星，其宏观孔隙度（大尺度结构性空腔）似乎极低[97]。

小行星的成分推断主要基于三类数据：反照率、表面光谱与密度。其中密度必须依赖对小行星卫星轨道的观测才能获得高精度结果。迄今所有具有卫星的小行星都被证实为碎石堆结构，即内部为岩石和金属碎块的松散集合体，其体积空隙甚至可高达 50\%。已调查的小行星最大直径可达 \(280\,\mathrm{km}\)，包括 121 Hermione（\(268\times186\times183\,\mathrm{km}\)）与 87 Sylvia（\(384\times262\times232\,\mathrm{km}\)）。超过 87 Sylvia 尺度的小行星极为罕见，且其中尚无被发现拥有卫星者。如此大尺度天体仍可能为碎石堆结构，且推定源自高度破坏性撞击事件，这一事实对太阳系形成模型具有重要影响：数值模拟显示，若碰撞双方为整体坚固体，则破坏概率与并合概率相当；而若双方为碎石堆结构，则并合概率显著提高。这意味着行星核心可能在形成早期即快速成长[98]。

\textbf{水}  

科学家推测，在形成月球的大碰撞之后，地球最初获得部分水的来源可能是来自小行星撞击[99]。2009 年，利用美国国家航空航天局（NASA）红外望远镜设施（Infrared Telescope Facility），科研团队在小行星 24 Themis 表面首次证实存在水冰，其表层似乎被冰层完全覆盖。由于该冰层正在发生升华，推测表面冰可能由地下冰储层持续补给。此外，还在其表面检测到有机化合物[100][101][99][102]。这些证据使“小行星向早期地球输送水”的假设更具可信性[99]。

2013 年 10 月，科研人员首次在太阳系外天体上探测到水，其目标为环绕白矮星 GD 61 的一颗小行星[103]。2014 年 1 月 22 日，欧洲航天局（ESA）宣布首次明确探测到位于小行星带最大天体 Ceres 的水汽[104]，该探测利用了赫歇尔空间天文台（Herschel Space Observatory）的远红外观测能力[105]。这一发现颇为意外，因为传统观点认为只有彗星才会喷发喷流与羽状气体；科研人员表示：“彗星与小行星之间的界限正变得越来越模糊”[105]。

研究发现，太阳风可与小行星表层矿物中的氧发生化学反应并生成水。据估算，以 Itokawa（S 型小行星）的样本为例，每立方米经辐射处理的岩石中可能可产生高达约 20 升的水[106][107]，其分析采用原子探针断层（atom probe tomography）技术。

1990 年在阿尔及利亚发现的陨石 Acfer 049 于 2019 年被证明具有超高孔隙岩性（ultraporous lithology，UPL），即由孔隙组成的结构，其孔隙可能曾被冰填充，后冰分离失去，因而推测 UPL 可能是“原初冰的化石”[108]。

\textbf{有机化合物}

小行星含有微量氨基酸及其他有机化合物，因此有观点认为，小行星撞击可能为早期地球提供了启动生命所需的化学原料，甚至可能直接携带生命至地球（即所谓“胚种论”，panspermia）[109][110]。2011 年 8 月，基于 NASA 对陨石样品研究的结果显示，DNA 与 RNA 的组成分子——腺嘌呤、鸟嘌呤及相关有机分子——可能在小行星与彗星中于外太空形成[111][112][113]。

2019 年 11 月，科学家首次在陨石中检测到糖类分子，包括核糖（ribose），这表明小行星可能具备合成生命基本必需生物分子的化学环境，从而支持 DNA 出现之前的“RNA 世界”假说，同时也为生命可能起源于外星并传播至地球的胚种论提供了额外证据[114][115][116]。
\subsection{分类}
小行星通常依据两类标准进行分类：其轨道特征，以及其反射光谱的特性。
\subsubsection{轨道分类}
许多小行星根据其轨道特征被归入不同的族群（groups）与家族（families）。除最广义的分类之外，通常以该组中最先被发现的成员来命名整个族群。族群通常是较为松散的动力学关联，而家族则更为紧密，被认为是过去某一时刻一颗大型母体小行星发生灾难性破碎后产生的碎片集合[117]。家族在主小行星带中更为常见，也更容易识别，但在木星特洛伊小行星中也已发现若干小规模家族[118]。主带家族最初由 Hirayama 清次于 1918 年识别，因此常以“平山家族”（Hirayama families）命名以示纪念。

大约 30–35\% 的主带小行星属于动力学家族，各家族被认为源自一次小行星间的过去碰撞。此外，冥族矮行星 Haumea 也被发现拥有一个相关的小行星家族。

部分小行星具有特殊的“马蹄形轨道”（horseshoe orbit），其轨道与地球或其他行星协轨。例如 3753 Cruithne 与 2002 AA29。此类轨道结构最早在土星卫星 Epimetheus 与 Janus 之间被发现。有时这些马蹄形轨道天体会暂时进入“准卫星”（quasi-satellite）状态，持续几十到几百年，然后再返回原先的马蹄形协轨状态。现已知地球与金星均存在准卫星。

若此类天体与地球、金星或假设与水星协轨，则它们属于 Aten 小行星的特殊类别。然而，它们也可能与外行星协轨。
\subsubsection{光谱分类} 
1975 年，Chapman、Morrison 与 Zellner 基于颜色、反照率及光谱形状建立了一套小行星分类系统[119]。这些性质被认为与小行星表面物质成分相关。最初的分类包含三类：C 型（暗色碳质小行星，占已知样本的约 75\%）、S 型（富硅石质小行星，占约 17\%）以及 U 型（无法归入 C 或 S 的未定类天体）。随着研究的深入，该体系被扩展并纳入更多类型，小行星光谱分类数量也持续增长。

目前最常使用的两类光谱分类体系分别为 Tholen 分类与 SMASS 分类。Tholen 分类由 David J. Tholen 于 1984 年提出，基于 1980 年代进行的八色光测量巡天数据，共划分出 14 类小行星[120]。2002 年发布的小行星主带光谱巡天（SMASS）在 Tholen 分类基础上做出修订，并扩展为 24 类。两种分类体系均包含三大主类：C 型、S 型与 X 型，其中 X 型主要包含金属富集类（如 M 型）小行星。此外还存在若干较小子类[121]。

不同光谱类型在已知小行星总体中所占比例并不必然反映其在真实总体中的占比，因为不同类型的小行星在观测可见性上存在偏差，导致统计结果存在观测偏倚。

\textbf{问题}

最初的光谱分类建立在对小行星成分的推断基础上[122]，但光谱类型与真实物质成分之间并非总能保持可靠对应，因此目前存在多套并行分类体系。这些差异造成了显著混淆。尽管不同光谱类别的小行星通常被认为由不同物质组成，但同一光谱类别内部的成员并不能保证具有相同或相似成分。
\subsubsection{活跃小行星}
\begin{figure}[ht]
\centering
\includegraphics[width=8cm]{./figures/54bec1528fc6a108.png}
\caption{由 OSIRIS\text{-}REx 拍摄的正在喷射颗粒物的 (101955) Bennu} \label{fig_Astero_15}
\end{figure}
活跃小行星是指轨道性质类似小行星，但外观呈现彗星特征的天体。也就是说，它们会出现彗发（coma）、尾迹（tail）或其他由质量损失所导致的可见特征（类似彗星），但其轨道仍位于木星轨道以内（类似小行星）[123][124]。该类天体最初由天文学家 David Jewitt 与 Henry Hsieh 于 2006 年命名为主带彗星（Main-Belt Comets, MBCs），但该名称暗示其必然具有类似彗星的冰质成分且仅存在于主带之内，而随着活跃小行星族群不断增长，这一推论已被证明并非总是成立[123][125][126]。

首个被发现的活跃小行星是 7968 Elst–Pizarro。它于 1979 年被视为小行星首次记录，但在 1996 年 Eric Elst 与 Guido Pizarro 发现其拥有尾迹后，被赋予彗星编号 133P/Elst–Pizarro[123][127]。另一个值得注意的例子是 311P/PanSTARRS：哈勃空间望远镜的观测显示其拥有六条类彗星尾迹[128]。这些尾迹被认为是由于碎石堆结构的小行星自转速度过快，导致表面物质被抛射形成的物质流[129]。
\begin{figure}[ht]
\centering
\includegraphics[width=8cm]{./figures/c7459379083513e3.png}
\caption{哈勃空间望远镜拍摄的 Dimorphos 及其在 DART 撞击后形成的尾迹} \label{fig_Astero_16}
\end{figure}
通过撞击小行星 Dimorphos，美国国家航空航天局（NASA）的“双小行星偏转测试”探测器（DART）使其成为一颗活跃小行星。科学家此前曾提出部分活跃小行星可能由撞击事件触发，但在此之前从未有人直接观测到小行星被激活的全过程。DART 任务在精确可控与严密观测的撞击条件下成功触发 Dimorphos 的活跃状态，从而首次实现对活跃小行星形成过程的细致研究[130][131]。观测显示，Dimorphos 在撞击后损失了约 \(1{,}000{,}000\,\mathrm{kg}\) 的物质[132]。撞击产生的尘埃羽流（dust plume）暂时性提高了 Didymos 系统的整体亮度，并形成了一条约 \(10{,}000\,\mathrm{km}\)（\(6{,}200\,\mathrm{mi}\)）长的尘埃尾迹，该尾迹持续存在数月之久[133][134][135]。
\subsubsection{暗彗星}  
暗彗星首次在 2024 年被识别，其轨道表现出非引力加速行为——这一特征通常见于彗星而非小行星——但却没有形成彗发或尾迹，因此在观测上与普通小行星无异。科研人员已确认其存在两类族群：其一为外族（Outer Family），成员直径大于 \(100\,\mathrm{m}\)，轨道可延伸至接近木星，与木星族彗星轨道相似；其二为内族（Inner Family），成员直径小于 \(50\,\mathrm{m}\)，反照率更低，轨道近似圆形[136]。除星际天体 'Oumuamua 外，所有暗彗星的轨道均位于主小行星带。值得注意的是，'Oumuamua 是首个在视像上呈小行星特征却表现出非引力加速的星际天体[137]。
\subsection{观测与探测}
在航天时代到来之前，主带小行星只能通过大型望远镜进行观测，其形状与地表细节长期未知。即使是当代最先进的地基望远镜与地球轨道上的哈勃空间望远镜，也仅能解析少数最大型小行星表面的有限细节。通过光变曲线（自转导致的亮度变化）与光谱特征可推断其形状与成分；通过掩星观测（小行星从恒星前方经过）可推算其尺寸。雷达成像则能获取较高质量的形状、轨道与自转参数信息，特别是针对近地小行星。航天器近距离飞掠探测可显著提升数据精度，而样品返回任务则可直接提供风化层（regolith）成分信息。
\subsubsection{地基观测}
\begin{figure}[ht]
\centering
\includegraphics[width=8cm]{./figures/1d3d6f9b77d6ee8b.png}
\caption{} \label{fig_Astero_17}
\end{figure}
\begin{figure}[ht]
\centering
\includegraphics[width=8cm]{./figures/78ef6c8773671d35.png}
\caption{} \label{fig_Astero_18}
\end{figure}
由于小行星体积较小且亮度微弱，地基观测（ground-based observations, GBO）所能获得的数据相当有限。通过地基光学望远镜可以测量其视星等，并进一步换算为绝对星等，从而得到其尺寸的粗略估计。地基观测还可获取光变曲线；若长时间累积数据，则可推算自转周期、（在部分情况下）自转极方向，以及小行星形状的粗略估计。可见光与近红外光谱观测可用于获取成分信息，并对目标进行分类。然而，由于此类观测仅能探测表层厚度为数微米的极薄外层，其限制明显[138]。正如行星科学家 Patrick~Michel 所述：

中红外至热红外观测与偏振测量，大概是仅有能够反映真实物理性质的数据。若只测量某一波长处的小行星热通量，可对对象尺寸作出估计，其不确定性通常低于基于可见光反射测量得到的估计。若两类数据同时获得，则可导出有效直径与几何反照率——后者定义为零相位角（即光源与观测者近似共线且光照来自观测者背后）时的表面亮度。此外，若同时获得多个热红外波长观测以及可见光亮度，则可进一步推断热学性质。多数已观测小行星的热惯量低于裸岩参考值，但高于月壤；这意味着其表面存在绝热作用的颗粒物层。并且观测表明存在某种趋势，可能与重力环境有关：体积较小（重力较弱）的天体表层覆盖较薄、颗粒较粗的风化层，而体积较大的天体则具有较厚、颗粒较细的风化层。然而，仅凭遥感观测仍难以获取该层材料的详细性质。此外，热惯量与表面粗糙度之间的关系并非线性明确，因此必须谨慎解释热惯量参量。

靠近地球的近地小行星可通过雷达成像获得更加细致的信息，例如表面是否存在陨石坑与巨石。相关观测由波多黎各的阿雷西博天文台（\SI{305}{m} 天线）以及加利福尼亚的戈德斯通深空站（\SI{70}{m} 天线）执行。雷达观测还可用于精确测定目标天体的轨道与自转动力学[138]。
\subsubsection{太空观测}
\begin{figure}[ht]
\centering
\includegraphics[width=8cm]{./figures/3ba5690de7058b16.png}
\caption{} \label{fig_Astero_19}
\end{figure}
无论是太空观测台还是地基观测台均开展了小行星巡天计划；其中太空观测预计可发现更多目标，因为不存在大气干扰，且可观测更大范围的天空。NEOWISE 观测了超过 100{,}000 颗主带小行星，斯皮策空间望远镜（Spitzer Space Telescope）观测了超过 700 颗近地小行星。上述观测为绝大多数目标提供了尺寸的粗略估计，但对表面性质（如风化层厚度与成分、安息角、内聚力、多孔度等）的细节信息仍较有限[138]。

哈勃空间望远镜（Hubble Space Telescope）也曾用于研究小行星，包括：追踪主带中发生碰撞的小行星[139][140]、观测小行星的解体过程[141]、观测具有六条类似彗尾结构的活跃小行星[142]、以及观测未来航天任务的目标小行星[143][144]。
\subsubsection{太空探测任务}  
根据 Patrick~Michel 的论述：

小行星内部结构仅能通过间接证据推断，例如：探测器测得的总体密度、对双星系统中天然卫星轨道的分析、以及由雅可夫斯基热效应导致的小行星轨道漂移。当探测器靠近小行星时，会受到其引力扰动，从而可以估算小行星的质量；小行星体积则需依据其形状模型进行估计。质量与体积可导出总体密度，其中不确定度通常由体积估计误差主导。通过将总体密度与推定陨石类比物的密度进行对比，可推断其内部孔隙率；观测表明暗色小行星往往更具高孔隙率（>40\%）而明亮小行星较低。但此类孔隙的具体性质仍不明确。

\textbf{专门探测任务}

首个被近距离成像的小行星为 951~Gaspra（1991~年），随后是 243~Ida 及其卫星 Dactyl（1993~年），均由前往木星途中经过的小行星并由伽利略号（Galileo）探测器拍摄。此外，还有数颗在任务途中被短暂访问的小行星，包括：9969~Braille（由“深空一号”Deep Space~1 于~1999~年探测）、5535~Annefrank（由“星尘号”Stardust 于~2002~年探测）、2867~Šteins 与~21~Lutetia（由“罗塞塔号”Rosetta 于~2008~年探测）、以及 4179~Toutatis（由中国探月工程“嫦娥二号”于~2012~年以~3.2~km（2~mi）的最近距离飞掠）。

首个真正面向小行星的专属探测器为 NASA 的 NEAR Shoemaker，其于~1997~年拍摄了 253~Mathilde，随后进入 433~Eros 的轨道并最终于~2001~年着陆其表面，成为历史上首个成功环绕并着陆小行星的航天器[145]。2005~年 9~月至~11~月，日本 隼鸟号（Hayabusa）探测器对 25143~Itokawa 进行了详细研究，并于~2010~年~6~月~13~日成功将样品带回地球，成为史上首个小行星样本返回任务。2007~年，NASA 发射了 黎明号
（Dawn）探测器，其先后环绕 4~Vesta（为期约一年）与矮行星 Ceres（为期约三年）。

2014~年，日本宇宙航空研究开发机构（JAXA）发射 隼鸟二号（Hayabusa2），于目标小行星 162173~Ryugu 轨道停驻超过一年，并成功采样，于~2020~年将样品送回地球。目前探测器正在执行延伸任务，预计将在~2031~年抵达新的目标。

NASA 于~2016~年发射 OSIRIS{\text -}REx，前往小行星 101955~Bennu 执行样本返回任务。该探测器于~2021~年离开小行星并携样品返回，样品于~2023~年~9~月送达地球。其后转入延伸任务并被命名为 OSIRIS{\text -}APEX，计划于~2029~年前往近地小行星 Apophis 探测。

2021~年，NASA 还发射了 DART（Double Asteroid Redirection Test）任务，用于测试可用于地球防御的偏转技术。DART 于~2022~年~9~月故意撞击双小行星系统 Didymos 的卫星 Dimorphos，目的是评估航天器撞击能否有效改变小行星轨道[146]。同年~10~月，NASA 宣布任务成功，确认 Dimorphos 绕 Didymos 的轨道周期缩短约~32~分钟[147]。

同于~2021~年发射的 NASA Lucy 是一项多小行星飞掠任务，目标为 7~颗不同类型的木星特洛伊小行星。尽管其尚未抵达首个主要目标 3548~Eurybates（预计~2027~年），但已先后飞掠主带小行星 152830~Dinkinesh 与 52246~Donaldjohanson[148][149]。

NASA 于~2023~年~10~月发射 Psyche 探测器，旨在研究同名的大型金属小行星，预计~2029~年抵达。

欧洲航天局（ESA）于~2024~年~10~月发射 Hera 探测器，用于研究 DART 撞击后的结果，包括撞击形成坑的形态与动量传递效率，以评估偏转技术效能。

中国国家航天局（CNSA）于~2025~年~5~月发射 天问二号[5]，将使用太阳能电推进访问共轨近地小行星 469219~Kamoʻoalewa 与活跃小行星 311P/PanSTARRS，并计划采集 Kamoʻoalewa 表层风化层样品[6]。



